% ---------------------------------------------------------------------------
% Preámbulo común del informe de calidad geométrica
% ---------------------------------------------------------------------------

% Codificación, idioma y tipografía
\usepackage[T1]{fontenc}
\usepackage{lmodern}
\usepackage[spanish]{babel}
\usepackage{amssymb}
\usepackage[left=2.5cm,right=2.5cm,top=2.5cm,bottom=2.5cm]{geometry}
\usepackage{microtype}
\usepackage{iftex} % \ifPDFTeX: permite usar opciones exclusivas de pdfTeX sin romper XeLaTeX

% Tablas, código y enlaces
\usepackage{xcolor}
\usepackage{booktabs}
\usepackage{tabularx}
\usepackage{enumitem}
\usepackage{float}
\usepackage{listings}
\usepackage{xstring}
\usepackage[hidelinks]{hyperref}

% Bibliografía según UNE-ISO 690:2013 (compilar con biber)
% Sistema autor-fecha (Harvard): \parencite{Autor, año} y \textcite{Autor (año)}
\usepackage{csquotes}
\usepackage[
	backend=biber,
	style=iso-authoryear,
	language=spanish,
	autolang=other,
	giveninits=true,
	uniquename=init,
	dateabbrev=false,
]{biblatex}
\addbibresource{bibliografia.bib}

% UNE-ISO 690:2013: «Disponible en: <URL> [Consulta: <fecha>]»
\DefineBibliographyStrings{spanish}{%
	urlfrom = {Disponible en},
	urlseen = {Consulta},
	online  = {en línea},
}
\DeclareFieldFormat{urldate}{\mkbibbrackets{\bibstring{urlseen}\addcolon\space#1}}

% La fecha de consulta se imprime detrás de la URL, como en la norma
\renewbibmacro*{urldate}{}
\renewbibmacro*{availability+access}{%
	\ifboolexpr{
		test {\iftoggle{bbx:doi}}
		and not test {\iffieldundef{doi}}
	}
	{\usebibmacro{from-doi}}
	{\ifboolexpr{
			test {\iftoggle{bbx:eprint}}
			and not test {\iffieldundef{eprint}}
		}
		{\usebibmacro{from-eprint}}
		{\usebibmacro{url}}}%
	\setunit{\addspace}%
	\printurldate}

\definecolor{codeblue}{RGB}{25,70,140}
\definecolor{codegray}{RGB}{245,245,245}

\lstdefinestyle{codeblock}{
	basicstyle=\ttfamily\small,
	keywordstyle=\color{codeblue}\bfseries,
	stringstyle=\color{red!60!black},
	breaklines=true,
	columns=fullflexible,
	frame=single,
	rulecolor=\color{gray!40},
	backgroundcolor=\color{codegray},
	showstringspaces=false,
	upquote=true,
	tabsize=2,
	captionpos=b
}
\lstset{style=codeblock}

% Los bloques de código se numeran como «Código» (no «Listing»)
\renewcommand{\lstlistingname}{Código}
\renewcommand{\lstlistlistingname}{Índice de códigos}

% Identificadores con caracteres especiales ($, _, etc.)
\ifPDFTeX
	\DisableLigatures[-]{encoding=T1,family=tt*}
\fi
\DeclareRobustCommand{\code}[1]{%
	\begingroup
	\StrSubstitute{\detokenize{#1}}{\detokenize{ñ}}{\~n}[\codeA]%
	\StrSubstitute{\codeA}{\string\$}{\textdollar}[\codeB]%
	\noexpandarg
	\expandafter\StrSubstitute\expandafter{\codeB}{'}{\textquotesingle}[\codeC]%
	\texttt{\codeC}%
	\endgroup
}
% En los marcadores del PDF (hiperenlaces) no puede ejecutarse xstring:
\pdfstringdefDisableCommands{\def\code#1{\detokenize{#1}}}
\newcommand{\seccion}[1]{Sección~\ref{#1}}

% Entradas del glosario de frontmatter/definiciones.tex
% \defterm{clave}{término} imprime únicamente la entrada de definición.
\newcommand{\defterm}[2]{\item[#2]}

\renewcommand{\arraystretch}{1.15}
